\chapter{KESIMPULAN DAN SARAN}

\section{Kesimpulan}

Penelitian ini berangkat dari permasalahan pelacakan \textit{invoice} dan \textit{Proof of Delivery} (POD) yang masih dilakukan secara semi-manual, sehingga status dokumen, bukti penerimaan, dan prioritas pengiriman belum terdokumentasi secara terintegrasi. Selain itu, penentuan prioritas pengiriman \textit{invoice} masih bergantung pada pengetahuan operasional yang belum sepenuhnya diformalkan, terutama terkait aturan \textit{cut-off}, jadwal penerimaan pelanggan, dan risiko keterlambatan administratif.

Pendekatan yang digunakan dalam penelitian ini menempatkan pengetahuan operasional sebagai dasar utama pembentukan mekanisme representasi pengetahuan dan rekomendasi prioritas. Melalui tahap \textit{Knowledge Acquisition}, pengetahuan diperoleh dari regulasi pelanggan, data historis \textit{invoice}, dan praktik administrasi yang relevan. Pengetahuan tersebut kemudian diproses melalui \textit{Knowledge Formalization} sehingga dapat direpresentasikan menjadi atribut operasional, \textit{Operational Labeling Guideline}, dan Rule Base yang eksplisit.

Hasil formalisasi pengetahuan digunakan dalam dua mekanisme representasi. Pertama, \textit{Rule-Based Representation} merepresentasikan pengetahuan operasional secara langsung melalui aturan IF-THEN R1 sampai R12. Kedua, \textit{Decision Tree Reconstruction} menggunakan atribut operasional dan \textit{Guideline-Based Ground Truth} untuk merekonstruksi pola keputusan. Dengan demikian, kedua mekanisme tidak ditempatkan sebagai pendekatan yang saling menggantikan, melainkan sebagai dua cara berbeda dalam merepresentasikan pengetahuan operasional yang sama.

Berdasarkan eksperimen akhir, \textit{Rule-Based Representation} menghasilkan rekomendasi prioritas yang konsisten dengan \textit{Operational Labeling Guideline} dan \textit{Guideline-Based Ground Truth}. Hal ini menunjukkan bahwa aturan yang diformalkan mampu merepresentasikan pedoman keputusan prioritas secara eksplisit dan dapat ditelusuri. Sementara itu, \textit{Decision Tree} C4.5 membentuk jalur keputusan yang sejalan dengan Rule Base, terutama melalui fitur operasional seperti jarak menuju \textit{cut-off}, kode hari penerimaan, jarak menuju hari penerimaan berikutnya, dan jadwal penerimaan. Temuan ini menunjukkan bahwa \textit{Decision Tree} C4.5 merekonstruksi pengetahuan operasional yang telah diformalkan, bukan menemukan pengetahuan operasional baru yang terpisah dari pedoman.

Kontribusi ilmiah utama dari penelitian ini adalah \textit{Operational Knowledge Formalization Framework}, yaitu kerangka yang menunjukkan bahwa pengetahuan operasional administrasi \textit{invoice} dapat diakuisisi, diformalkan, direpresentasikan dalam bentuk aturan eksplisit, dan direkonstruksi kembali melalui \textit{Decision Tree Reconstruction}. Implikasi ilmiahnya adalah bahwa evaluasi tidak harus dibaca sebagai pemeringkatan metode, tetapi sebagai \textit{Knowledge Representation Consistency} antara representasi eksplisit dan rekonstruksi pengetahuan. Implikasi praktisnya adalah rekomendasi prioritas dapat ditempatkan dalam sistem pelacakan \textit{invoice} dan POD agar keputusan administratif lebih tertelusur.

\section{Keterbatasan}

Penelitian ini memiliki beberapa keterbatasan yang perlu diperhatikan. Pertama, dataset final yang digunakan masih berukuran relatif kecil, sehingga variasi pola operasional yang tercakup belum dapat merepresentasikan seluruh kemungkinan kondisi administrasi \textit{invoice}. Ukuran dataset tersebut memadai untuk menunjukkan konsistensi antara pedoman operasional, Rule Base, dan rekonstruksi pengetahuan, tetapi belum cukup untuk menarik generalisasi yang luas.

Kedua, dataset berasal dari konteks satu perusahaan. Oleh karena itu, aturan \textit{cut-off}, jadwal penerimaan, dan pola prioritas yang digunakan sangat dipengaruhi oleh karakteristik pelanggan dan proses administrasi pada lingkungan tersebut. Hasil penelitian tidak dapat langsung digeneralisasi ke perusahaan lain tanpa penyesuaian terhadap regulasi pelanggan dan proses operasional yang berbeda.

Ketiga, \textit{Guideline-Based Ground Truth} pada eksperimen dibentuk melalui \textit{Guideline-Based Labeling} berdasarkan \textit{Operational Labeling Guideline}. Pendekatan ini sesuai dengan tujuan penelitian, yaitu mengevaluasi konsistensi antara representasi eksplisit dan rekonstruksi pengetahuan dari pengetahuan operasional yang sama. Namun, label tersebut tidak dapat ditafsirkan sebagai validasi eksternal yang selalu identik dengan keputusan manusia pada seluruh situasi nyata.

Keempat, validasi eksternal terhadap data keputusan organisasi nyata belum dilakukan secara luas. Dengan demikian, kemampuan kerangka untuk mempertahankan konsistensi representasi pada periode waktu lain, pelanggan baru, atau perubahan regulasi administratif masih perlu diuji lebih lanjut. Perubahan aturan pelanggan juga dapat memengaruhi relevansi Rule Base dan struktur pohon keputusan yang telah terbentuk.

\section{Saran Penelitian Selanjutnya}

Penelitian selanjutnya dapat memperluas cakupan data dengan menggunakan dataset yang lebih besar, berasal dari periode waktu yang lebih panjang, dan memuat keputusan organisasi nyata sebagai bahan validasi. Penambahan data tersebut dapat membantu menguji stabilitas konsistensi representasi dan meningkatkan pemahaman terhadap variasi kondisi administrasi \textit{invoice}.

Validasi pada lebih dari satu perusahaan juga perlu dilakukan agar generalisasi kerangka dapat diuji pada konteks operasional yang berbeda. Setiap perusahaan dapat memiliki aturan \textit{cut-off}, jadwal penerimaan, dan prosedur verifikasi yang berbeda, sehingga pengujian lintas organisasi dapat memperkuat validitas eksternal penelitian.

Selain itu, penelitian selanjutnya dapat melibatkan lebih banyak sumber pengetahuan operasional, termasuk beberapa staf administrasi atau pihak yang memahami regulasi pelanggan. Pendekatan ini dapat memperkaya \textit{Operational Labeling Guideline} dan mengurangi ketergantungan pada satu sumber interpretasi operasional.

Dari sisi evaluasi organisasi, penelitian berikutnya dapat membandingkan keluaran rekomendasi prioritas dengan keputusan aktual staf administrasi pada proses kerja nyata. Validasi ini penting untuk menilai apakah \textit{Operational Labeling Guideline}, Rule Base, dan rekonstruksi pohon keputusan tetap sesuai dengan praktik keputusan organisasi ketika digunakan di luar dataset penelitian.

Penelitian berikutnya juga dapat menelaah proses pembaruan pedoman ketika regulasi pelanggan berubah. Fokus pengembangan sebaiknya tetap diarahkan pada validasi keputusan organisasi, keterlacakan pengetahuan, dan konsistensi representasi, bukan pada peningkatan algoritma semata.

Pada tahap implementasi sistem, penelitian selanjutnya dapat mengintegrasikan rekomendasi prioritas dengan sistem ERP atau sistem administrasi perusahaan secara langsung. Integrasi tersebut dapat mendukung rekomendasi operasional secara real-time, pembaruan status \textit{invoice}, penyimpanan POD, serta pemantauan perubahan regulasi pelanggan. Selain itu, mekanisme pembaruan pengetahuan dapat dipertimbangkan apabila data operasional baru terus tersedia dan proses bisnis mengalami perubahan secara berkala.
